\section{Class PSPACE --- TQBFHard Verifier}
%%% generated by the blueprint pipeline from TCSlib.Complexity.ClassPSPACE.TQBFHardVerifier
\subsection*{Class PSPACE --- TQBFHard Verifier: Introduction}

Declarations of the Lean module \texttt{TCSlib.Complexity.ClassPSPACE.TQBFHardVerifier}.

\subsection*{Class PSPACE --- TQBFHard Verifier: Results}

\begin{theorem}[Polynomial-time closure under truncation]
\lean{Complexity.TQBFHard.pt_drop}
\label{Complexity.TQBFHard.pt_drop}
\leanok
\uses{Complexity.PolyTimeComputable, Complexity.PolyTimeComputable.comp, Complexity.PolyTimeComputable.pairEncode, Complexity.PrefixByLength.drop, Complexity.polyTimeComputable_dropPrefixByLength}
\difficulty{3}
Let $o\colon\{0,1\}^*\to\mathbb N$ and $S\colon\{0,1\}^*\to\{0,1\}^*$ be polynomial-time
computable. Then the function sending each bit string $y$ to the string obtained by
deleting the first $o(y)$ bits of $S(y)$ is polynomial-time computable.
\end{theorem}

\begin{theorem}[Polynomial-time prefix truncation]
\lean{Complexity.TQBFHard.pt_take}
\label{Complexity.TQBFHard.pt_take}
\leanok
\uses{Complexity.PolyTimeComputable, Complexity.PolyTimeComputable.comp, Complexity.PolyTimeComputable.pairEncode, Complexity.PrefixByLength.take, Complexity.polyTimeComputable_takePrefixByLength}
\difficulty{3}
Let $o$ assign a natural number to each bit string, and let $S$ be a function from bit
strings to bit strings. If the unary string of length $o(y)$ is polynomial-time
computable as a function of $y$, and $S$ is polynomial-time computable, then the
function sending $y$ to the prefix of $S(y)$ of length $o(y)$ is polynomial-time
computable.
\end{theorem}

\begin{theorem}[Polynomial-time unary replication]
\lean{Complexity.TQBFHard.pt_rep_mul}
\label{Complexity.TQBFHard.pt_rep_mul}
\leanok
\uses{Complexity.PolyTimeComputable, Complexity.PolyTimeComputable.append, Complexity.PolyTimeComputable.comp, Complexity.polyTimeComputable_const, Complexity.polyTimeComputable_unary}
\difficulty{4}
For every $k\in\mathbb N$ and every polynomial-time computable function
$g\colon\{0,1\}^*\to\{0,1\}^*$, the function sending $y$ to the string of $k|g(y)|$ ones
is polynomial-time computable.
\end{theorem}

\begin{definition}[Keep the true entries]
\lean{Complexity.TQBFHard.keepTrue}
\label{Complexity.TQBFHard.keepTrue}
\leanok
\uses{Complexity.transduce}
For a list of Boolean values, define its true-entry list by deleting every false entry
and retaining each true entry, in its original order.
\end{definition}

\begin{theorem}[Filtering true bits]
\lean{Complexity.TQBFHard.keepTrue_eq_nil_iff}
\label{Complexity.TQBFHard.keepTrue_eq_nil_iff}
\leanok
\uses{Complexity.TQBFHard.keepTrue, Complexity.transduce}
\difficulty{4}
For every finite Boolean list $l$, removing all false entries from $l$ yields the empty
list if and only if every entry of $l$ is false.
\end{theorem}

\begin{theorem}[Equality via bitwise XOR]
\lean{Complexity.TQBFHard.eq_iff_xor}
\label{Complexity.TQBFHard.eq_iff_xor}
\leanok
\difficulty{4}
For bit strings $u$ and $v$, they are equal if and only if they have the same length and
every bit in their pointwise XOR is false.
\end{theorem}

\begin{theorem}[Polynomial-time equality test]
\lean{Complexity.TQBFHard.pt_eq}
\label{Complexity.TQBFHard.pt_eq}
\leanok
\uses{Complexity.PolyTimeComputable, Complexity.PolyTimeComputable.comp, Complexity.PolyTimeComputable.pairEncode, Complexity.TQBFHard.eq_iff_xor, Complexity.TQBFHard.keepTrue_eq_nil_iff, Complexity.pairFstD_pairEncode, Complexity.pairSndD_pairEncode, Complexity.polyTimeComputable_and, Complexity.polyTimeComputable_isNil, Complexity.polyTimeComputable_lenEq, Complexity.polyTimeComputable_transduce, Complexity.polyTimeComputable_xorD, Complexity.xorD_pairEncode}
\difficulty{5}
Let $f,g\colon\{0,1\}^*\to\{0,1\}^*$ be polynomial-time computable. The function that
outputs the one-bit string $[1]$ when $f(y)=g(y)$ and $[0]$ otherwise is polynomial-time
computable.
\end{theorem}

\begin{theorem}[Polynomial-time initial-configuration indicator]
\lean{Complexity.TQBFHard.pt_indicator}
\label{Complexity.TQBFHard.pt_indicator}
\leanok
\uses{Complexity.LOGSPACE, Complexity.LOGSPACE_subset_P, Complexity.PolyTimeComputable, Complexity.mem_P_iff, Turing.MultiTapeTM.indicator}
\difficulty{4}
Let $L$ be a language of bit strings decidable in logarithmic space. For each bit string
$z$, let $x$ and $r$ be the first and second components returned by the total pair
projections of $z$. Form the encoding whose first component is $x$ and whose second
component pairs the length-$vwid(A,a,|x|)$ prefix of $r$ beginning at offset
$vwid(A,a,|x|)$ with the binary expansion of $vrad(d,c,|x|)$. The function outputting
the one-bit string whose bit is the indicator of membership of this encoding in the
initial-configuration language for $N$ is polynomial-time computable.
\end{theorem}

\begin{definition}[Verification radius]
\lean{Complexity.TQBFHard.vrad}
\label{Complexity.TQBFHard.vrad}
\leanok
Given a nondeterministic Turing machine $N$ with Boolean input alphabet and
natural-number parameters $d$ and $c$, define the verification radius at $n$ to be
$2d(n+1)^c$.
\end{definition}

\begin{definition}[Verification width]
\lean{Complexity.TQBFHard.vwid}
\label{Complexity.TQBFHard.vwid}
\leanok
Given a nondeterministic Turing machine $N$ with Boolean input alphabet and
natural-number parameters $A$ and $a$, define the verification width at $n$ to be
$A(n+1)^a$.
\end{definition}

\begin{definition}[Verification code length]
\lean{Complexity.TQBFHard.vlen}
\label{Complexity.TQBFHard.vlen}
\leanok
\uses{Complexity.TQBFHard.codeLen, Complexity.TQBFHard.vrad}
Given a nondeterministic Turing machine $N$ with Boolean input alphabet and
natural-number parameters $d$ and $c$, define the verification code length at $n$ to be
the code length of $N$ on input length $n$ with radius $2d(n+1)^c$.
\end{definition}

\begin{definition}[Extracted code segment]
\lean{Complexity.TQBFHard.cpAt}
\label{Complexity.TQBFHard.cpAt}
\leanok
\uses{Complexity.TQBFHard.vlen, Complexity.pairFstD, Complexity.pairSndD}
Given a nondeterministic Turing machine $N$ with Boolean input alphabet, natural-number
parameters $d$ and $c$, a Boolean list $y$, and a natural number $\mathrm{off}$, let
$y_1$ and $y_2$ be the first and second components obtained by decoding $y$. The
extracted segment is the portion of $y_2$ after its first $\mathrm{off}$ entries,
truncated to the verification code length for input length equal to the length of $y_1$.
\end{definition}

\begin{definition}[Extracted configuration block]
\lean{Complexity.TQBFHard.blkAt}
\label{Complexity.TQBFHard.blkAt}
\leanok
\uses{Complexity.TQBFHard.vwid, Complexity.pairFstD, Complexity.pairSndD}
For a Boolean list $y$ and a natural number $r$, write $n$ for the length of the first
component of the decoded pair represented by $y$. The block at offset $r$ is the segment
of length $vwid(A,a,n)$ beginning at position $r$ in the second component of that
decoded pair.
\end{definition}

\begin{definition}[Initial configuration predicate]
\lean{Complexity.TQBFHard.initB}
\label{Complexity.TQBFHard.initB}
\leanok
\uses{Complexity.Codec.initLang, Complexity.TQBFHard.cpAt, Complexity.TQBFHard.vrad, Complexity.TQBFHard.vwid, Complexity.pairFstD, Turing.MultiTapeTM.indicator, Turing.pairEncode}
For a nondeterministic Turing machine $N$ with Boolean tape alphabet, natural numbers
$d,c,A,a$, and a Boolean list $y$, write $x$ for the first component of the decoded pair
represented by $y$. The initial-configuration predicate is true exactly when the
machine's initial-language indicator accepts the encoding of $x$, the configuration
specified by $N,d,c,y$ and width $vwid(A,a,|x|)$, and the binary representation of
$vrad(d,c,|x|)$.
\end{definition}

\begin{definition}[Accepting configuration predicate]
\lean{Complexity.TQBFHard.accB}
\label{Complexity.TQBFHard.accB}
\leanok
\uses{Complexity.Codec.acceptLang, Complexity.TQBFHard.cpAt, Complexity.TQBFHard.vrad, Complexity.TQBFHard.vwid, Complexity.pairFstD, Turing.MultiTapeTM.indicator, Turing.pairEncode}
For a nondeterministic Turing machine $N$ with Boolean tape alphabet, natural numbers
$d,c,A,a$, and a Boolean list $y$, write $x$ for the first component of the decoded pair
represented by $y$. The accepting-configuration predicate is true exactly when the
machine's accepting-language indicator accepts the encoding of $x$, the configuration
specified by $N,d,c,y$ and width $2vwid(A,a,|x|)$, and the binary representation of
$vrad(d,c,|x|)$.
\end{definition}

\begin{definition}[Base configuration predicate]
\lean{Complexity.TQBFHard.baseB}
\label{Complexity.TQBFHard.baseB}
\leanok
\uses{Complexity.Codec.stepLang, Complexity.TQBFHard.blkAt, Complexity.TQBFHard.cpAt, Complexity.TQBFHard.vrad, Complexity.TQBFHard.vwid, Complexity.pairFstD, Turing.MultiTapeTM.indicator, Turing.pairEncode}
For a nondeterministic Turing machine $N$ with Boolean tape alphabet, natural numbers
$d,c,A,a$, and a Boolean list $y$, put $x$ for the first component of the decoded pair
represented by $y$ and $w=vwid(A,a,|x|)$. The base predicate is true if either the
blocks of length $w$ in the second component of the decoded pair at offsets $3w^2+w$ and
$3w^2+2w$ are equal, or the machine's step-language indicator accepts the encoding of
$x$, the two configurations specified by $N,d,c,y$ at those respective offsets, and the
binary representation of $vrad(d,c,|x|)$.
\end{definition}

\begin{definition}[Selected witness block]
\lean{Complexity.TQBFHard.winBlk}
\label{Complexity.TQBFHard.winBlk}
\leanok
\uses{Complexity.pairFstD, Complexity.pairSndD}
For a Boolean list $q$ and a natural number $k$, the witness block at index $k$ is the
segment of the second component of the decoded pair represented by $q$, starting at
position $k$ times the length of the second component of the decoded first component,
and having that same length.
\end{definition}

\begin{definition}[Selection consistency predicate]
\lean{Complexity.TQBFHard.selT}
\label{Complexity.TQBFHard.selT}
\leanok
\uses{Complexity.TQBFHard.winBlk}
For a Boolean list $q$, the selection predicate is true exactly when either the witness
blocks at indices $4$ and $1$ are equal and those at indices $5$ and $3$ are equal, or
the blocks at indices $4$ and $3$ are equal and those at indices $5$ and $2$ are equal.
\end{definition}

\begin{definition}[Level query encoding]
\lean{Complexity.TQBFHard.lvlQ}
\label{Complexity.TQBFHard.lvlQ}
\leanok
\uses{Complexity.TQBFHard.vwid, Complexity.pairFstD, Complexity.pairSndD, Turing.pairEncode}
For a Boolean list $y$ and a natural number $j$, write $x$ for the first component of
the decoded pair represented by $y$. The level query encoding pairs the encoding of $x$
together with a string of $vwid(A,a,|x|)$ true bits, and the suffix of the second
component of the decoded pair beginning at position $3vwid(A,a,|x|)j$.
\end{definition}

\begin{definition}[Existential selection failure predicate]
\lean{Complexity.TQBFHard.orB}
\label{Complexity.TQBFHard.orB}
\leanok
\uses{Complexity.TQBFHard.lvlQ, Complexity.TQBFHard.selT, Complexity.TQBFHard.vwid, Complexity.pairFstD}
For a nondeterministic Turing machine $N$ with Boolean tape alphabet, natural numbers
$d,c,A,a$, and a Boolean list $y$, the predicate is true exactly when there is some
natural number $j$ strictly less than $vwid(A,a,|x|)$, where $x$ is the first component
of the decoded pair represented by $y$, for which the selection predicate is false on
the level query encoding of $y$ at level $j$.
\end{definition}

\begin{definition}[Combined Boolean predicate]
\lean{Complexity.TQBFHard.predB}
\label{Complexity.TQBFHard.predB}
\leanok
\uses{Complexity.TQBFHard.accB, Complexity.TQBFHard.baseB, Complexity.TQBFHard.initB, Complexity.TQBFHard.orB}
For a Boolean-valued nondeterministic Turing machine $N$ and natural numbers $d,c,A,a$,
define a predicate on Boolean lists $y$ that is true exactly when
$	ext{initB}(N,d,c,A,a,y)$ and $	ext{accB}(N,d,c,A,a,y)$ are both true, and at least one
of $	ext{baseB}(N,d,c,A,a,y)$ and $	ext{orB}(A,a,y)$ is true.
\end{definition}

\begin{theorem}[Polynomial-time unary encoding of a width bound]
\lean{Complexity.TQBFHard.pt_vwid}
\label{Complexity.TQBFHard.pt_vwid}
\leanok
\uses{Complexity.PolyTimeComputable, Complexity.PolyTimeComputable.comp, Complexity.TQBFHard.vwid, Complexity.pairFstD, Complexity.polyTimeComputable_pairFstD, Complexity.polyTimeComputable_polyUnary}
\difficulty{1}
Let $A,a\in\mathbb N$. The function sending a bit string $y$ to a unary string of length
$A(|x|+1)^a$, where $x$ is the decoded first component of $y$ and is empty if decoding
fails, is polynomial-time computable.
\end{theorem}

\begin{theorem}[Polynomial-time unary encoding of a quadratic width expression]
\lean{Complexity.TQBFHard.pt_vwid_lin}
\label{Complexity.TQBFHard.pt_vwid_lin}
\leanok
\uses{Complexity.PolyTimeComputable, Complexity.PolyTimeComputable.append, Complexity.PolyTimeComputable.comp, Complexity.TQBFHard.vwid, Complexity.pairFstD, Complexity.polyTimeComputable_pairFstD, Complexity.polyTimeComputable_polyUnary}
\difficulty{4}
Let $A,a,k,l\in\mathbb N$. The function sending a bit string $y$ to a unary string of
length $kV^2+lV$ is polynomial-time computable, where $V=A(|x|+1)^a$ and $x$ is the
decoded first component of $y$ (empty if decoding fails).
\end{theorem}

\begin{theorem}[Polynomial-Time Computability of the Radius Bits]
\lean{Complexity.TQBFHard.pt_bits_vrad}
\label{Complexity.TQBFHard.pt_bits_vrad}
\leanok
\uses{Complexity.PolyTimeComputable, Complexity.PolyTimeComputable.comp, Complexity.TQBFHard.vrad, Complexity.pairFstD, Complexity.polyTimeComputable_of_linear, Complexity.polyTimeComputable_pairFstD, Complexity.polyTimeComputable_polyUnary, Turing.FinTM.computesFunInTime_lengthBits}
\difficulty{3}
For every nondeterministic finite Turing machine $N$ and natural numbers $d,c,A,a$, the
function sending a Boolean word $y$ to the binary expansion, in the specified list
representation, of $2d(|x|+1)^c$ is polynomial-time computable, where $x$ is the first
component extracted from $y$.
\end{theorem}

\begin{theorem}[Polynomial-time unary encoding of a configuration length]
\lean{Complexity.TQBFHard.pt_vlen}
\label{Complexity.TQBFHard.pt_vlen}
\leanok
\uses{Complexity.Codec.W, Complexity.Codec.wP, Complexity.PolyTimeComputable, Complexity.PolyTimeComputable.append, Complexity.PolyTimeComputable.comp, Complexity.TQBFHard.codeLen, Complexity.TQBFHard.vlen, Complexity.TQBFHard.vrad, Complexity.pairFstD, Complexity.polyTimeComputable_const, Complexity.polyTimeComputable_of_linear, Complexity.polyTimeComputable_pairFstD, Complexity.polyTimeComputable_polyUnary, Complexity.polyTimeComputable_unary, Turing.FinTM.computesFunInTime_lengthBits}
\difficulty{4}
Let $N$ be a finite nondeterministic Turing machine over the binary alphabet, and let
$d,c$ be natural numbers. For each bit string $y$, form its decoded first component $x$
(taking the empty string when decoding fails). The function that outputs a unary string
of length $V_{N,d,c}(|x|)$ is polynomial-time computable, where
$V_{N,d,c}(n)=W_N(4d(n+1)^c+1)+|\operatorname{bits}(n+1)|+1$ and $W_N=|H_N|+3k_N+1$;
here $k_N$ is the number of work tapes, $H_N$ is the finite set of
state-and-output-summary pairs, and $\operatorname{bits}(m)$ is the binary digit list of
$m$.
\end{theorem}

\begin{theorem}[Polynomial-Time Computability of a Configuration Slice]
\lean{Complexity.TQBFHard.pt_cpAt}
\label{Complexity.TQBFHard.pt_cpAt}
\leanok
\uses{Complexity.PolyTimeComputable, Complexity.TQBFHard.cpAt, Complexity.TQBFHard.pt_drop, Complexity.TQBFHard.pt_take, Complexity.TQBFHard.pt_vlen, Complexity.polyTimeComputable_pairSndD}
\difficulty{1}
Fix a nondeterministic finite Turing machine $N$ and natural numbers $d,c,A,a$. Let $o$
assign a natural number $o(y)$ to each Boolean word $y$, and suppose the function
sending $y$ to a word of $o(y)$ true bits is polynomial-time computable. Then the
function sending $y$ to the slice of the second component extracted from $y$, beginning
at offset $o(y)$ and having length equal to the configuration-code length determined by
the length of the first component of $y$, is polynomial-time computable.
\end{theorem}

\begin{theorem}[Polynomial-Time Computability of a Block Slice]
\lean{Complexity.TQBFHard.pt_blkAt}
\label{Complexity.TQBFHard.pt_blkAt}
\leanok
\uses{Complexity.PolyTimeComputable, Complexity.TQBFHard.blkAt, Complexity.TQBFHard.pt_drop, Complexity.TQBFHard.pt_take, Complexity.TQBFHard.pt_vwid, Complexity.polyTimeComputable_pairSndD}
\difficulty{1}
Fix a nondeterministic finite Turing machine $N$ and natural numbers $d,c,A,a$. Let $o$
assign a natural number $o(y)$ to each Boolean word $y$, and suppose the function
sending $y$ to a word of $o(y)$ true bits is polynomial-time computable. Then the
function sending $y$ to the slice of the second component extracted from $y$, beginning
at offset $o(y)$ and having length $A(|x|+1)^a$, is polynomial-time computable, where
$x$ is the first component extracted from $y$.
\end{theorem}

\begin{theorem}[Polynomial-time initial-configuration test]
\lean{Complexity.TQBFHard.pt_initB}
\label{Complexity.TQBFHard.pt_initB}
\leanok
\uses{Complexity.Codec.initLang_mem, Complexity.PolyTimeComputable, Complexity.PolyTimeComputable.comp, Complexity.PolyTimeComputable.pairEncode, Complexity.TQBFHard.initB, Complexity.TQBFHard.pt_bits_vrad, Complexity.TQBFHard.pt_cpAt, Complexity.TQBFHard.pt_indicator, Complexity.TQBFHard.pt_vwid, Complexity.polyTimeComputable_pairFstD}
\difficulty{3}
For every finite nondeterministic Turing machine $N$ over the binary alphabet and all
natural numbers $d,c,A,a$, the function sending a bit string $y$ to the one-bit string
containing the value of the following Boolean test is polynomial-time computable: take
the first pair projection $x$ of $y$, form the encoding pairing $x$ with a pair
consisting of the length-$vwid(A,a,|x|)$ prefix of the second projection of $y$ starting
at offset $vwid(A,a,|x|)$ and the binary expansion of $vrad(d,c,|x|)$, and test whether
that encoding belongs to the initial-configuration language for $N$.
\end{theorem}

\begin{theorem}[Polynomial-Time Computability of an Acceptance Bit]
\lean{Complexity.TQBFHard.pt_accB}
\label{Complexity.TQBFHard.pt_accB}
\leanok
\uses{Complexity.Codec.acceptLang_mem, Complexity.PolyTimeComputable, Complexity.PolyTimeComputable.comp, Complexity.PolyTimeComputable.pairEncode, Complexity.TQBFHard.accB, Complexity.TQBFHard.pt_bits_vrad, Complexity.TQBFHard.pt_cpAt, Complexity.TQBFHard.pt_indicator, Complexity.TQBFHard.pt_vwid_lin, Complexity.TQBFHard.vwid, Complexity.pairFstD, Complexity.polyTimeComputable_pairFstD}
\difficulty{3}
For every nondeterministic finite Turing machine $N$ and natural numbers $d,c,A,a$, the
function sending a Boolean word $y$ to the one-bit word whose bit is true exactly when
the encoded input and configuration data extracted from $y$ form an accepting
configuration in the language recognized by $N$ is polynomial-time computable.
\end{theorem}

\begin{theorem}[Polynomial-Time Computability of a Base Test]
\lean{Complexity.TQBFHard.pt_baseB}
\label{Complexity.TQBFHard.pt_baseB}
\leanok
\uses{Complexity.Codec.stepLang_mem, Complexity.PolyTimeComputable, Complexity.PolyTimeComputable.comp, Complexity.PolyTimeComputable.pairEncode, Complexity.TQBFHard.baseB, Complexity.TQBFHard.pt_bits_vrad, Complexity.TQBFHard.pt_blkAt, Complexity.TQBFHard.pt_cpAt, Complexity.TQBFHard.pt_eq, Complexity.TQBFHard.pt_indicator, Complexity.TQBFHard.pt_vwid_lin, Complexity.polyTimeComputable_or, Complexity.polyTimeComputable_pairFstD}
\difficulty{4}
For every nondeterministic finite Turing machine $N$ and natural numbers $d,c,A,a$, the
function sending a Boolean word $y$ to the one-bit word whose bit is true when either
the two extracted blocks at offsets $3w^2+w$ and $3w^2+2w$ agree, or the encoded data
certify a valid step of $N$, is polynomial-time computable; here $w=A(|x|+1)^a$ and $x$
is the first component extracted from $y$.
\end{theorem}

\begin{theorem}[Polynomial-time selector test]
\lean{Complexity.TQBFHard.pt_selT}
\label{Complexity.TQBFHard.pt_selT}
\leanok
\uses{Complexity.PolyTimeComputable, Complexity.PolyTimeComputable.comp, Complexity.TQBFHard.pt_drop, Complexity.TQBFHard.pt_eq, Complexity.TQBFHard.pt_rep_mul, Complexity.TQBFHard.pt_take, Complexity.TQBFHard.selT, Complexity.TQBFHard.winBlk, Complexity.pairFstD, Complexity.pairSndD, Complexity.polyTimeComputable_and, Complexity.polyTimeComputable_or, Complexity.polyTimeComputable_pairFstD, Complexity.polyTimeComputable_pairSndD, Complexity.polyTimeComputable_unary}
\difficulty{4}
For every choice of the parameters $N,d,c,A,a$, the function sending a bit string $q$ to
the one-bit string containing the selector value is polynomial-time computable. The
selector is true exactly when either the fourth and first blocks of $q$ agree and the
fifth and third blocks agree, or the fourth and third blocks agree and the fifth and
second blocks agree, where block $i$ is the length-$|b|$ segment of the second pair
projection of the first pair projection of $q$ beginning at offset $i|b|$, and $b$ is
the second pair projection of the first pair projection of $q$.
\end{theorem}

\begin{definition}[Next encoded word]
\lean{Complexity.TQBFHard.nxtW}
\label{Complexity.TQBFHard.nxtW}
\leanok
\uses{Complexity.pairFstD, Complexity.pairSndD, Turing.pairEncode}
For a Boolean-valued nondeterministic Turing machine $N$ and natural numbers $d,c,A,a$,
define a transformation on Boolean lists $q$ as follows: pair-encode the first decoded
component of $q$ with the second decoded component after deleting its first $3$ times
the length of the second decoded component of the first decoded component of $q$
entries.
\end{definition}

\begin{theorem}[Polynomial-time bounded disjunction test]
\lean{Complexity.TQBFHard.pt_orB}
\label{Complexity.TQBFHard.pt_orB}
\leanok
\uses{Complexity.BlockLoop.polyTimeComputable_orLoop, Complexity.PolyTimeComputable, Complexity.PolyTimeComputable.comp, Complexity.PolyTimeComputable.pairEncode, Complexity.TQBFHard.lvlQ, Complexity.TQBFHard.nxtW, Complexity.TQBFHard.orB, Complexity.TQBFHard.pt_drop, Complexity.TQBFHard.pt_rep_mul, Complexity.TQBFHard.pt_selT, Complexity.TQBFHard.pt_vwid, Complexity.TQBFHard.selT, Complexity.TQBFHard.vwid, Complexity.length_pair_components_le, Complexity.pairFstD, Complexity.pairFstD_pairEncode, Complexity.pairSndD, Complexity.pairSndD_pairEncode, Complexity.polyTimeComputable_not, Complexity.polyTimeComputable_pairFstD, Complexity.polyTimeComputable_pairSndD, Turing.length_pairEncode, Turing.pairEncode}
\difficulty{6}
For every choice of the parameters $N,d,c,A,a$, the function sending a bit string $y$ to
the one-bit string containing the Boolean value of the following test is polynomial-time
computable. Let $x$ be the first pair projection of $y$, and set $m=A(|x|+1)^a$. The
test is true exactly when there is some $j<m$ for which the selector evaluated on the
encoding pairing the encoding of $x$ with the length-$m$ string of ones and the suffix
of the second projection of $y$ after its first $3mj$ bits is false.
\end{theorem}

\begin{theorem}[Polynomial-time computability of the predicate bit]
\lean{Complexity.TQBFHard.polyTimeComputable_predB}
\label{Complexity.TQBFHard.polyTimeComputable_predB}
\leanok
\uses{Complexity.PolyTimeComputable, Complexity.TQBFHard.predB, Complexity.TQBFHard.pt_accB, Complexity.TQBFHard.pt_baseB, Complexity.TQBFHard.pt_initB, Complexity.TQBFHard.pt_orB, Complexity.polyTimeComputable_and, Complexity.polyTimeComputable_or}
\difficulty{1}
Let $N$ be a finite nondeterministic Turing machine over the Boolean alphabet, and let
$d,c,A,a$ be natural numbers. The function that maps each finite Boolean string $y$ to
the one-bit string whose bit is true exactly when the initial-configuration check and
accepting-configuration check for $N$ on the encoded input both hold, and either the
base check or the disjunction of level checks holds, is computable by a deterministic
finite Turing machine in time bounded by a polynomial in the input length.
\end{theorem}
